\documentclass[10pt,a4paper]{article}
\usepackage[margin=0.52in]{geometry}
\usepackage{amsmath,amssymb}
\usepackage{booktabs}
\usepackage{tabularx}
\usepackage{xcolor}
\usepackage{colortbl}
\usepackage{hyperref}

\definecolor{navy}{HTML}{0D1117}
\definecolor{slate}{HTML}{1B2838}
\definecolor{rowtint}{HTML}{F1F5F9}

\setlength{\parindent}{0pt}
\setlength{\parskip}{4pt}
\renewcommand{\arraystretch}{1.22}

\begin{document}

% Header Banner — Replace placeholders with your own title, name, and Student ID
\noindent\colorbox{navy}{\parbox{\dimexpr\textwidth-2\fboxsep\relax}{%
\vspace{8pt}
\hspace{8pt}{\color{white}\Large\bfseries YOUR PROJECT TITLE: Six-Scene Educational Storyboard}\\[4pt]
\hspace{8pt}{\color[HTML]{8AEAD8}\normalsize Your Focused Mathematical Subtitle \quad|\quad M3 First-Year Undergraduate Mathematics}\\[6pt]
\hspace{8pt}{\color[HTML]{CBD5E1}\small \textbf{Student:} Your Name \quad\textbullet\quad \textbf{Student ID:} YOUR\_ID \quad\textbullet\quad \textbf{Format:} 1920$\times$1080, 30 fps, 16:9 MP4 \quad\textbullet\quad \textbf{Target:} 5:00}
\vspace{8pt}
}}

\vspace{8pt}

\section*{\color{slate}1. Learning Objective \& Mathematical Scope}
\textbf{Learning Objective:} By the end of this video, the viewer will be able to explain \textit{[insert your focused concept]}, verify \textit{[insert your worked example]}, and interpret \textit{[insert the geometric/structural meaning]}.

\vspace{4pt}
\noindent\colorbox{rowtint}{\parbox{\dimexpr\textwidth-2\fboxsep\relax}{%
\vspace{4pt}
\begin{tabularx}{\linewidth}{@{}l X@{}}
\textbf{Core Definition / Rule:} & \textit{Insert your main formula, e.g., $A\vec{v} = \lambda\vec{v}$ or $\nabla f(x,y) = (\partial f/\partial x, \partial f/\partial y)$} \\[4pt]
\textbf{Worked Example:} & \textit{Insert your concrete numerical example and verified final result} \\[4pt]
\textbf{Geometric Meaning:} & \textit{State what the final numbers mean visually}
\end{tabularx}
\vspace{4pt}
}}

\section*{\color{slate}2. Hybrid Visual System \& Audio Specification}
\begin{tabularx}{\textwidth}{@{}l X l X@{}}
\toprule
\textbf{Canvas} & Deep charcoal/navy (\texttt{\#0D1117}) & \textbf{Primary Highlight} & Cyan (\texttt{\#53D8FB}) \\
\textbf{Reference Layer} & Subdued blue/grey (\texttt{\#526477}) & \textbf{Secondary Highlight} & Amber (\texttt{\#FFB84D}) \\
\textbf{Active / Transformed} & Restrained bright teal (\texttt{\#55D6BE}) & \textbf{Tracked Point} & Coral (\texttt{\#FF718D}) \\
\textbf{Narration Voice} & Gemini TTS (\texttt{Kore}, $\sim$500--560 words) & \textbf{Music / Assets} & No background music (or credit track here) \\
\bottomrule
\end{tabularx}

\section*{\color{slate}3. Six-Scene Production Storyboard}

\small
\begin{tabularx}{\textwidth}{@{}p{2.0cm} p{1.45cm} X X@{}}
\toprule
\rowcolor{slate}
{\color{white}\textbf{Scene}} & {\color{white}\textbf{Timing}} & {\color{white}\textbf{Visual Action \& On-Screen Formulae}} & {\color{white}\textbf{Narration Focus \& Transition}} \\
\midrule

\textbf{1. Hook / Context} &
0:00--0:30\newline ($\sim$30\,s) &
Opening title frame (title, concept, your name \& ID). Present an intriguing visual puzzle or geometric question. &
Motivate why the concept matters before introducing heavy formulas. Smooth fade to Scene 2. \\
\midrule

\rowcolor{rowtint}
\textbf{2. Concept Intro} &
0:30--1:05\newline ($\sim$35\,s) &
Introduce formal notation, objects, and definitions beside a clean geometric diagram. &
Define every symbol clearly in natural spoken English before relying on it. \\
\midrule

\textbf{3. Visual Intuition} &
1:05--2:10\newline ($\sim$65\,s) &
Animate the core visual idea (grid, vector, curve, tangent, surface, or level set) step by step. &
Build geometric understanding first so the viewer sees \textit{why} the math works. \\
\midrule

\rowcolor{rowtint}
\textbf{4. Worked Example} &
2:10--3:15\newline ($\sim$65\,s) &
Show one carefully chosen numerical example with every intermediate calculation displayed cleanly. &
Connect each arithmetic step directly to what moves or changes on screen. \\
\midrule

\textbf{5. Interpretation} &
3:15--4:25\newline ($\sim$70\,s) &
Connect the final numerical result back to geometry, scaling, direction, or optimality. &
Explain assumptions, edge cases, and what the final value means visually. \\
\midrule

\rowcolor{rowtint}
\textbf{6. Summary \& Credits} &
4:25--5:00\newline ($\sim$35\,s) &
Display three concise takeaway cards, followed by references, music credit, and AI disclosure. &
Reinforce the learning objective clearly and close with transparent credits. \\
\bottomrule
\end{tabularx}

\end{document}

